\begin{textsample}{POS Dim 1 – pe – Score 28.00 – pe/pe\_ef\_40.txt}  \label{ex:f1_pos_190}
3.6.2 ARTE COMO CONHECIMENTO A \textbf{abordagem} mais \textbf{contemporânea} do ensino de Arte está relacionada ao desenvolvimento cognitivo, \textbf{que} envolve flexibilidade, fluência, elaboração e nos impõe pensar de maneira diferente o ensino de Arte na educação escolar.
Permite deslocar as questões de “ como se ensina Arte ” para “ como se aprende Arte ”.
A partir dessa mudança de paradigmas, vários estudos vieram a ser desenvolvidos, buscando explicar o processo de \textbf{ensino-aprendizagem} dos conhecimentos artísticos.
Essa \textbf{abordagem} visa à construção de \textbf{um} currículo para o ensino de Arte \textbf{que} interligue a leitura da Arte, sua \textbf{contextualização} e o fazer artístico produzido pela \textbf{humanidade}, \textbf{que} engloba a produção internacional, nacional, regional, local e dos próprios estudantes.
Ressignificando esses paradigmas de \textbf{ensino-aprendizagem} dos conhecimentos artísticos e num contexto de luta dos profissionais do ensino de Arte \textbf{aqui} no Brasil, foi promulgada a LDB, de nº 9.394/96, \textbf{que} garantiu a obrigatoriedade do ensino de Arte para toda a educação básica.
A Lei assegurou \textbf{que} o Ensino de Arte deverá promover o desenvolvimento cultural dos educandos, fundamentado na concepção da Arte como conhecimento.
Ressalta -se ainda \textbf{que} o ensino de Arte como conhecimento está \textbf{baseado} nos estudos culturais, no interculturalismo, na \textbf{interdisciplinaridade} e na aprendizagem dos conhecimentos artísticos, a partir da \textbf{inter-relação} entre o fazer, o ler e o contextualizar Arte.
A Arte Contemporânea é caracterizada pelo rompimento de barreiras entre o visual, o gestual e o sonoro.
O happening, a performance ( em suas múltiplas \textbf{faces} ), a body art, a arte sociológica e ambiental, o conceitualismo, as artes híbridas e a própria videoarte são algumas das manifestações artísticas \textbf{que} comprovam \textbf{uma} tendência \textbf{atual} para o inter-relacionamento de diversas linguagens no ensino de Arte ( BARBOSA, 2002b ).
A concepção de Arte como conhecimento busca a valorização tanto do produto artístico como dos processos \textbf{desencadeados} no ensino de Arte, percorrendo o caminho contrário ao da concepção de ensino como técnica – \textbf{que} valoriza o produto artístico em detrimento do processo – e da concepção de ensino de Arte como expressão – \textbf{que} supervaloriza o processo, dando pouca importância ao produto estético, ao mesmo tempo, sem desconsiderar a técnica e os recursos tecnológicos \textbf{que} vêm como instrumento para a produção artística.
Na educação escolar, \textbf{enquanto} componente curricular, a Arte abrange todas as linguagens artísticas, conforme a Lei nº 13.278/2016, e compõe -se de quatro grandes campos distintos de conhecimento : Artes Visuais, Dança, Música e Teatro, incluídos nos currículos das etapas de educação infantil e ensino fundamental.
A nova lei altera a Lei de Diretrizes e Bases da Educação Nacional ( LDB 9.394/1996 ), estabelecendo prazo de cinco anos para \textbf{que} os sistemas de ensino promovam a formação de professores para implantar esses componentes curriculares em todos os níveis da educação básica no ensino infantil, fundamental e médio ( BRASIL, 2016 ).
Cada \textbf{um} deles tem suas especificidades, \textbf{exigindo} profissionais com formação específica.
Os quatro campos têm \textbf{interfaces} entre si.
Contemporaneamente, também se deve levar em consideração outras formas híbridas de arte, tais como performance, webarte e multimídia, por exemplo, \textbf{que} se apresentam em diversos espaços e em diferentes nuances ( PERNAMBUCO, 2013 ).
Não podemos \textbf{perder} de vista a colaboração do ensino-aprendizagem de Arte para o desenvolvimento integral dos educandos dos diversos níveis de ensino, buscando abarcar as \textbf{inúmeras} possibilidades de criação e fruição artísticas frente às tecnologias disponíveis no mundo \textbf{contemporâneo}.
Apresentamos, brevemente, cada \textbf{um} dos quatro grandes campos, sua importância e relação com as habilidades e objetos de conhecimento possíveis no currículo em nosso estado.
O campo das Artes Visuais é extenso e múltiplo.
Sua matéria-prima é a imagem.
Contempla expressões bidimensionais como desenho, fotografia, gravura e pintura, por exemplo, bem como tridimensionais, como a escultura e a instalação, além de formas híbridas de expressões, a exemplo da videoperformance e da arte eletrônica, dentre outras.
Temáticas como as linhas, as formas, a composição, a cor, a luz e sombra, o volume, a perspectiva, as texturas, as diversas imagens veiculadas nas diferentes mídias, os movimentos artísticos, as correntes estéticas, bem como histórias e obras de artistas, as manifestações \textbf{populares}, as culturas de tradição, as histórias de vida são possibilidades a serem (re)construídas, (re)criadas, (re)inventadas pelo/a professor/a no espaço escolar e/ou fora dele.
Os princípios e \textbf{fundamentos} da arte da Dança podem ser trabalhados em sala de aula, tomando -se como base a diversidade das danças cênicas e \textbf{populares}, tais como o frevo, o maracatu, o caboclinho, o cavalo marinho, entre outras manifestações locais, regionais, nacionais ou internacionais, bem como o balé, a dança \textbf{moderna} e a dança \textbf{contemporânea}.
A Dança possui como matéria-prima o corpo.
As danças devem ser abordadas em seu contexto estético, histórico, social e cultural, sem \textbf{perder} de vista o incentivo à criação da própria arte de movimento do estudante.
Tais princípios devem ser abordados, para além dos conteúdos relativos a repertórios específicos, por meio de técnicas e metodologias não diretivas, \textbf{que} introduzam a improvisação como \textbf{estímulo} à criação e à composição coreográfica.
O campo da Música possui como matéria-prima o som e \textbf{baseia} -se numa \textbf{abordagem} sociocultural da educação musical \textbf{que} \textbf{revela} a importância de considerar as músicas das diferentes culturas, inclusive da escola e de seu entorno, e o trabalho com essas músicas como possibilidades diferenciadas de organização sonora e de meios de ampliação da experiência e discurso musicais dos estudantes.
O universo do Teatro é construído pelo \textbf{ator}, público, texto e espaço.
O Teatro, como \textbf{um} dos campos da Arte, traz a possibilidade de nos fazer refletir sobre as nossas subjetividades, sobre a vida e o lugar em \textbf{que} \textbf{vivemos}.
O ensino de Teatro na Escola possui \textbf{uma} trajetória histórica e socioepistemológica marcada por diferentes \textbf{abordagens} : as \textbf{pedagogias} do Jogo e Teatro, da Dramaturgia, da \textbf{Contextualização} Histórica, da Formação do \textbf{Ator}, do Teatro Épico, da Mediação Sociocultural, da Performance, das Manifestações Tradicionais, do Teatro de Formas Animadas e da Encenação.
Por fim, o ensino de Arte, no ensino fundamental – anos iniciais e finais, está comprometido com o desenvolvimento integral e cultural dos educandos.

% matched lemmas: abordagem, aqui, ator, atual, basear, contemporâneo, contextualização, desencadear, enquanto, ensino-aprendizagem, estímulo, exigir, face, fundamento, humanidade, inter-relação, interdisciplinaridade, interface, inúmero, moderno, pedagogia, perder, popular, que, revelar, um, vivar
\end{textsample}
